%%%%%%%% IEEE TPAMI SURVEY LATEX TEMPLATE %%%%%%%%%%%%%%%%%

\documentclass[journal]{IEEEtran}

\usepackage{cite}
\usepackage{amsmath,amssymb,amsfonts}
\usepackage{graphicx}
\usepackage{booktabs}
\usepackage{hyperref}
\usepackage{xcolor}
\usepackage{multirow}
\usepackage{tabularx}
\usepackage[capitalize,noabbrev]{cleveref}

\graphicspath{{../figures/}} % To reference your generated figures, name the PNGs directly. DO NOT CHANGE THIS.

\begin{filecontents}{references.bib}
@article{lecun2015deep,
  title={Deep learning},
  author={LeCun, Yann and Bengio, Yoshua and Hinton, Geoffrey},
  journal={Nature},
  volume={521},
  number={7553},
  pages={436--444},
  year={2015},
  publisher={Nature Publishing Group}
}
\end{filecontents}

\title{Computer Vision for Construction Drawing Understanding: A Comprehensive Survey}

\author{Zelin~Wan
\thanks{Z. Wan is with Bobyard Inc., San Francisco, CA, USA.}
}

\begin{document}

\maketitle

%%%%%%%%%ABSTRACT%%%%%%%%%
\begin{abstract}
ABSTRACT HERE
\end{abstract}
%%%%%%%%%ABSTRACT%%%%%%%%%

\begin{IEEEkeywords}
construction drawings, technical drawing understanding, computer vision, deep learning, survey
\end{IEEEkeywords}

%%%%%%%%%INTRODUCTION%%%%%%%%%
\section{Introduction}
\label{sec:introduction}
% CONTENT HERE
%%%%%%%%%INTRODUCTION%%%%%%%%%

%%%%%%%%%BACKGROUND%%%%%%%%%
\section{Background and Taxonomy}
\label{sec:background}
% CONTENT HERE
%%%%%%%%%BACKGROUND%%%%%%%%%

%%%%%%%%%STATISTICAL_OVERVIEW%%%%%%%%%
\section{Statistical Overview of the Field}
\label{sec:statistical_overview}
% CONTENT HERE

% EXAMPLE FIGURE: REPLACE AND ADD YOUR OWN FIGURES / CAPTIONS
\begin{figure}[t]
\centering
\includegraphics[width=\columnwidth]{example-image-a}
\caption{PLEASE FILL IN CAPTION HERE}
\label{fig:statistical_overview}
\end{figure}
%%%%%%%%%STATISTICAL_OVERVIEW%%%%%%%%%

%%%%%%%%%SYMBOL_DETECTION%%%%%%%%%
\section{Symbol Detection and Recognition}
\label{sec:symbol_detection}
% CONTENT HERE
%%%%%%%%%SYMBOL_DETECTION%%%%%%%%%

%%%%%%%%%LAYOUT_ANALYSIS%%%%%%%%%
\section{Layout Analysis and Segmentation}
\label{sec:layout_analysis}
% CONTENT HERE
%%%%%%%%%LAYOUT_ANALYSIS%%%%%%%%%

%%%%%%%%%TEXT_RECOGNITION%%%%%%%%%
\section{Text Detection and Recognition in Drawings}
\label{sec:text_recognition}
% CONTENT HERE
%%%%%%%%%TEXT_RECOGNITION%%%%%%%%%

%%%%%%%%%TOPOLOGY%%%%%%%%%
\section{Topology and Graph Extraction}
\label{sec:topology}
% CONTENT HERE
%%%%%%%%%TOPOLOGY%%%%%%%%%

%%%%%%%%%FLOOR_PLANS%%%%%%%%%
\section{Architectural Floor Plans}
\label{sec:floor_plans}
% CONTENT HERE
%%%%%%%%%FLOOR_PLANS%%%%%%%%%

%%%%%%%%%ENGINEERING_DIAGRAMS%%%%%%%%%
\section{P\&ID and Engineering Diagram Analysis}
\label{sec:engineering_diagrams}
% CONTENT HERE
%%%%%%%%%ENGINEERING_DIAGRAMS%%%%%%%%%

%%%%%%%%%FOUNDATION_MODELS%%%%%%%%%
\section{Foundation Model Adaptation for Drawings}
\label{sec:foundation_models}
% CONTENT HERE
%%%%%%%%%FOUNDATION_MODELS%%%%%%%%%

%%%%%%%%%VLM%%%%%%%%%
\section{Vision-Language Models for Construction Documents}
\label{sec:vlm}
% CONTENT HERE
%%%%%%%%%VLM%%%%%%%%%

%%%%%%%%%HIGH_RESOLUTION%%%%%%%%%
\section{High-Resolution and Multi-Scale Processing}
\label{sec:high_resolution}
% CONTENT HERE
%%%%%%%%%HIGH_RESOLUTION%%%%%%%%%

%%%%%%%%%DATASETS%%%%%%%%%
\section{Datasets and Benchmarks}
\label{sec:datasets}
% CONTENT HERE
%%%%%%%%%DATASETS%%%%%%%%%

%%%%%%%%%CHALLENGES%%%%%%%%%
\section{Open Challenges and Future Directions}
\label{sec:challenges}
% CONTENT HERE
%%%%%%%%%CHALLENGES%%%%%%%%%

%%%%%%%%%CONCLUSION%%%%%%%%%
\section{Conclusion}
\label{sec:conclusion}
% CONTENT HERE
%%%%%%%%%CONCLUSION%%%%%%%%%

\bibliographystyle{IEEEtran}
\bibliography{references}

\end{document}
